\topsection[sec:intro]{Introduction}

\subsection{Context}
Machine-learning systems are assembled more than they are written. A model that reaches production descends
from a dataset collected by one team, a pretrained model downloaded from a public hub, a fine-tuning run that may
happen on someone else's infrastructure, and a packaging step. Each hand-over is a point where an artifact can be
replaced or manipulated.

The threat is concrete. In December 2022 a dependency-confusion attack replaced \code{torchtriton}, a dependency
of the PyTorch nightly builds, with a package that stole SSH keys and files \cite{pytorch2022}. In February 2024
JFrog found about one hundred models on Hugging Face whose pickle files ran code as soon as they were loaded
\cite{jfrog2024}. Poisoning is cheap: web-scale training sets can be poisoned for tens of dollars
\cite{carlini2024}, and in 2025 about 250 poisoned documents were enough to plant a backdoor in language models
of 600 million to 13 billion parameters \cite{souly2025}. Regulation follows. Article 15(5) of the EU AI Act
requires, for high-risk systems, measures against attacks that manipulate the training data (data poisoning) or
the pretrained components used in training (model poisoning) \cite{aiact2024}.

Software supply-chain security offers mature tools for part of the problem. SLSA, in-toto and Sigstore give a
build a verifiable provenance \cite{slsa11,intoto2019,sigstore2022}, and OpenSSF Model Signing applies signing to
model directories \cite{oms2025}. Two gaps remain. First, these tools attest one build or one file. Nothing
links the dataset to the base model, the base model to the fine-tuned model and that model to what was deployed,
with a rule on who may sign which step and a revocation that follows the links. Second, a signature proves that
bytes are what a key holder attested. It does not prove that the bytes behave well. For source code the
difference is tolerable, because code can be reviewed. For a file of weights it is not.

\subsection{Problem statement}
The project answers the question set by the brief of Projet Applicatif 2:

\begin{keybox}
How can an enterprise cryptographically verify that an AI model has not been tampered with at any stage of its
supply chain, from data collection through the acquisition of a pretrained base model to fine-tuning and
deployment?
\end{keybox}

Two sub-questions structure the work. The first is about \emph{integrity}: can we detect any modification,
substitution or re-linking of an artifact after a trusted party attested it? The second is about the
\emph{trust boundary}: what remains unprotected, in particular malicious content introduced before attestation,
and what evidence other than cryptography can narrow that gap?

\subsection{Contributions}
\begin{enumerate}
\item \textbf{A working attestation system.} \sys\ signs every artifact of a four-stage pipeline, links each
  one to its parents by hash, verifies an artifact and its whole ancestry with ten checks, keeps keys and
  revocations in a signed, pinned, append-only trust log, and gates deployment on the result. It has a command
  line, a REST API, a dashboard, a Docker Compose stack and two continuous-integration workflows
  (Sections~\ref{sec:design} and~\ref{sec:impl}).
\item \textbf{A reproducible evaluation.} 48 attack and control variants in seven experiments, 168 automated
  tests and a set of cost measurements, all regenerated by scripts in the repository
  (Section~\ref{sec:eval}).
\item \textbf{A measured limit.} Every manipulation made after signing is detected. A dataset poisoned before
  signing verifies and deploys with a backdoor success rate of 77.2\,\%, an adaptive trigger evades the content
  checks, and a backdoored model from a trusted publisher verifies with 100\,\%. The status is the same for ten
  other random seeds (Sections~\ref{sec:eval} and~\ref{sec:discussion}).
\item \textbf{Valorisation.} A six-page paper in IEEE conference format, a public website on which the signed
  model is verified and run in the reader's browser, and a public repository (Section~\ref{sec:valo}).
\end{enumerate}

\subsection{How this report follows the brief}
The brief defines six phases. Table~\ref{tab:phases} gives, for each phase, the section of this report and the
project documents that hold the detailed deliverable. Appendix~\ref{app:trace} lists every requirement of the
brief with its evidence.

\begin{table}[H]
\caption{The six phases of the brief, the sections of this report and the detailed deliverables in the repository.}
\label{tab:phases}
\tablefont
\begin{tabularx}{\textwidth}{P{4.5cm}P{2.3cm}L}
\toprule
\headrow\hd{Phase of the brief} & \hd{Section} & \hd{Detailed deliverables (folder \code{docs/} unless stated)} \\
\midrule
1\enskip Scoping and problem definition & \ref{sec:scoping} & Scoping note (\code{01}) \\
2\enskip State of the art and critical analysis & \ref{sec:soa} & State-of-the-art synthesis (\code{02}), comparative table (\code{03}) \\
3\enskip Solution design & \ref{sec:threat}, \ref{sec:design} & Threat model (\code{04}), architecture (\code{05}), functional and technical descriptions (\code{06}, \code{07}), attack scenarios (\code{08}) \\
4\enskip Implementation and experimentation & \ref{sec:impl} & Source code, 168 tests, Docker and CI files, experimentation journal (\code{09}) \\
5\enskip Testing, evaluation and critical analysis & \ref{sec:eval}, \ref{sec:discussion} & \code{results/} (logs, tables, figures), attack scenarios (\code{08}) \\
6\enskip Reporting and valorisation & \ref{sec:valo}, \ref{sec:project} & This report, the paper (\code{paper/}), the conference proposal (\code{13}), the website \\
\bottomrule
\end{tabularx}
\end{table}
